Le tissu economique tunisien repose massivement sur des milliers de petites
et moyennes entreprises (PME) de service : cliniques, cabinets medicaux,
laboratoires d'analyses, salons de coiffure, garages automobiles, cabinets
juridiques et professionnels independants. Toutes ces structures partagent
un mode de fonctionnement quotidien remarquablement uniforme : les rendez-vous
sont notes dans un agenda papier, les communications avec la clientele
transitent par WhatsApp depuis le telephone personnel du proprietaire ou
d'un secretaire, et la coordination entre les differents canaux d'information
(agenda, conversations, facturation, dossier client) repose entierement sur
la memoire et l'attention de l'operateur humain.

Cette organisation, qui fonctionne tant que le volume reste modeste, devient
rapidement le goulet d'etranglement de l'activite : les messages recus en
dehors des heures de travail restent sans reponse, le proprietaire est
constamment interrompu par des questions repetitives, les langues utilisees
par la clientele (arabe dialectal tunisien, francais, anglais, souvent
melangees dans une meme phrase) imposent une charge cognitive supplementaire,
et aucune visibilite analytique n'est possible sur les flux de l'entreprise.

L'evolution recente des technologies d'intelligence artificielle, en
particulier la disponibilite de modeles de langage grand public (Large
Language Models, ou LLM) capables de comprendre l'arabe dialectal et de
produire des reponses structurees, ouvre une opportunite : celle d'automatiser
le travail de receptionniste lui-meme, en l'integrant directement a l'interieur
du canal que la clientele utilise deja, c'est-a-dire WhatsApp.

Le present rapport presente le projet \textbf{Nova~AI}, une plateforme
multi-tenant de Software-as-a-Service (SaaS) qui realise concretement cette
opportunite. Nova~AI combine une intelligence artificielle conversationnelle
basee sur Anthropic Claude, l'API Cloud officielle de WhatsApp Business, une
base de donnees Postgres multi-locataire avec une securite au niveau des
lignes (Row-Level Security), un tableau de bord temps reel pour le
proprietaire, et une application web progressive (PWA) destinee aux clients
finaux.

Le projet a ete structure autour de la methodologie \textbf{Agile Scrum},
adaptee aux contraintes d'une equipe d'une seule personne et aux iterations
courtes imposees par le rythme du semestre. Quatre sprints successifs ont
permis de produire un systeme complet, deploye en production, capable de
recevoir et traiter des bookings reels via WhatsApp en arabe dialectal,
francais et anglais.

\section*{Plan du rapport}
Ce rapport est organise comme suit. Le \textbf{Chapitre I} presente le
cadre general du projet : l'organisme d'accueil, l'analyse du contexte
tunisien, l'etude de l'existant, la problematique et la solution proposee,
ainsi que la methodologie Scrum adoptee. Le \textbf{Chapitre II} (Sprint 0)
detaille la phase preparatoire : identification des acteurs, analyse des
besoins fonctionnels et non fonctionnels, choix techniques, diagrammes de
conception et planification des sprints. Les \textbf{Chapitres III, IV et V}
documentent les trois sprints de realisation : le sprint des fondations, le
sprint des surfaces operationnelles, et le sprint de la couche
d'intelligence. Enfin, la \textbf{Conclusion generale} fait le bilan du
travail accompli, identifie les limites du systeme actuel, et trace les
perspectives d'evolution future.
